\documentclass[11pt]{article}
\usepackage[a4paper,margin=25mm]{geometry}
\usepackage{amsmath,amssymb,amsthm,mathtools}
\usepackage[hidelinks]{hyperref}
\allowdisplaybreaks[2]
\setlength{\emergencystretch}{2em}
\newtheorem{theorem}{Theorem}
\newtheorem{lemma}[theorem]{Lemma}
\newtheorem{proposition}[theorem]{Proposition}
\title{Vertex extensions and tight cores in Tuza's conjecture}
\author{Leo Gao\\[-1pt]{\small AI-assisted research}}
\date{November 26, 2026 \quad$\cdot$\quad Research draft}
\begin{document}
\maketitle
\begin{abstract}
We prove Tuza's inequality for graphs obtained by adjoining a vertex to a graph whose triangle components are copies of $K_4$ or $K_5$. A second extension argument shows that, after removal of the edges of any triangle of an edge-minimal counterexample, the remaining graph is tight, has packing number one smaller, and contains an exceptional triangle component. Exact restoration formulas give an equivalent packing-augmentation condition for the full conjecture. Small-graph checks test the local conditions and verify a counterexample to unrestricted simultaneous wedge avoidance.
\end{abstract}

\section{Definitions}

Write $\nu(G)$ for the maximum size of an edge-disjoint triangle packing and $\tau(G)$ for the minimum number of edges meeting all triangles of a finite simple graph $G$. Tuza conjectured that $\tau(G)\le2\nu(G)$~\cite{Tuza}. A graph is \emph{tight} if $\tau(G)=2\nu(G)$.

The triangle components of $G$ are the unions of triangles in the connected components of the graph whose vertices are triangles of $G$ and whose adjacencies indicate a shared edge. Say that $G$ has \emph{standard triangle structure} when each triangle component is a $K_4$ or a $K_5$. Edges in no triangle are allowed. The two graph parameters add over triangle components because their edge sets are disjoint. For $K_4$ the pair $(\nu,\tau)$ is $(1,2)$. For $K_5$ it is $(2,4)$ by Mantel's theorem and a balanced bipartition for the cover, and because three packed triangles would use nine edges and have even degree at every vertex.

The results below concern extensions and minimal counterexamples. The line-graph case appears in the companion manuscript.

\section{Arbitrary vertex-extension theorem}

\textbf{Theorem.} Let \(H\) be any finite simple graph whose triangle components are copies of \(K_4\) or \(K_5\). Construct \(G\) by adding a new vertex with an arbitrary neighborhood in \(H\). Then

\[
\boxed{\tau(G)\leq2\nu(G).}
\]

There is no restriction on the degree of the added vertex, the number of components, or the presence of additional triangle-free edges in \(H\).

This provides a new exact extension mechanism beyond the triangle-deletion approach developed previously.

The new vertex may connect arbitrarily to multiple triangle components.

\paragraph{Proof.}

Write

\[
k=\nu(H).
\]

Because the triangle components have disjoint edge sets, and because \(K_4\) and \(K_5\) both attain equality in Tuza's inequality,

\[
\tau(H)=2k.
\]

Let \(v\) be the new vertex, and define its link graph

\[
L=H[N_G(v)].
\]

Choose an inclusion-maximal matching \(M\) in \(L\), subject to the restriction that at most one matching edge belongs to each \(K_4\) triangle component.

No additional restriction is needed for \(K_5\) components.

\textbf{Packing construction.}

A \(K_4\) contains a maximum triangle packing avoiding any specified edge.

A \(K_5\) contains a maximum triangle packing avoiding any two disjoint specified edges. Indeed, if those edges are \(ab\) and \(cd\), and the fifth vertex is \(e\), the triangles \(ace\) and \(bde\) provide the required packing.

Consequently, we can choose a maximum triangle packing of \(H\) avoiding every edge of \(M\).

For each \(xy\in M\), add the triangle \(vxy\).

These additional triangles are pairwise edge-disjoint, because \(M\) is a matching.

Hence

\[
\nu(G)\geq k+|M|.
\]

\textbf{Cover construction.}

Delete every edge joining \(v\) to an endpoint of \(M\). This costs exactly \(2|M|\) edges.

Every remaining triangle through \(v\) corresponds to an edge of \(L\) whose endpoints are both unmatched.

By the maximality of \(M\), such an edge must belong to a \(K_4\) component already represented by an edge of \(M\).

Inside that component, the only possible edge with both endpoints unmatched is the edge opposite the selected matching edge.

Choose the minimum triangle cover of this \(K_4\) to contain both opposite edges. This is a valid cover of size two.

Choose minimum covers independently in the remaining components.

The resulting cover of \(G\) has size at most

\[
2k+2|M|.
\]

Therefore,

\[
\boxed{
\tau(G)\leq2k+2|M|\leq2\nu(G).
}
\]

\[
\square
\]

\section{Restoring two incident edges}

Suppose \(J\) has standard triangle structure and \(H\) is obtained from \(J\) by adding two distinct incident edges.

Then it is impossible to have simultaneously

\[
\nu(H)=\nu(J)+1
\]

and

\[
\tau(H)=2\nu(H).
\]

\paragraph{Proof.}

Write the two added edges as \(wu,wv\), and put

\[
k=\nu(J).
\]

Assume for contradiction that

\[
\nu(H)=k+1,\qquad \tau(H)=2k+2.
\]

There must be a triangle using \(wu\) but not \(wv\), and another using \(wv\) but not \(wu\).

Otherwise, a minimum cover of \(J\), together with one additional edge, would cover \(H\).

Write these triangles as

\[
wux,\qquad wvy.
\]

\textbf{Case I: \(x\neq y\).}

The two triangles are edge-disjoint.

We show that a maximum triangle packing of \(J\) can avoid all four edges

\[
wx,\ ux,\ wy,\ vy.
\]

No complete triangle component of \(J\) contains both outside edges of either new triangle: doing so would require one of the missing edges \(wu,wv\).

Moreover, whenever a component contains two of the four edges, those edges must be incident.

A disjoint pair consisting of one edge through $w$ and one of $ux,vy$ cannot lie in a single component: that component would contain both endpoints of one of the absent edges $wu,wv$. If $ux$ and $vy$ belong to the same component, it contains $u,v,x,y$ and hence $xy$. The triangle $wxy$ then puts $w$ in the component, again forcing an absent edge.

Thus each component must avoid at most two incident edges.

Both cliques admit maximum packings avoiding two incident edges $ab,ac$. For $K_4$ use the triangle on the other three vertices. For $K_5$ with remaining vertices $d,e$, use $ade$ and $bce$.

Combining these packings and adding the two new triangles gives

\[
\nu(H)\geq k+2,
\]

a contradiction.

\textbf{Case II: \(x=y\) for every choice of extension triangles.}

All triangles using exactly one new edge therefore share the same outside vertex \(x\).

They can all be covered by the single edge \(wx\).

The only possible additional triangle is \(wuv\).

If this triangle exists, then \(uvx\) is a triangle of \(J\). Choose a minimum cover of its complete triangle component containing \(uv\).

Together with \(wx\), this covers all triangles of \(H\) using at most

\[
2k+1
\]

edges.

Again, this contradicts the assumed equality.

\[
\square
\]

\section{A necessary condition for a minimal counterexample}

This theorem produces a stronger necessary condition for a minimal counterexample than our previous open-wedge reduction.

\textbf{Theorem.} If \(G\) is an edge-minimal counterexample to Tuza's conjecture, with

\[
\nu(G)=k,
\]

then, for every triangle \(T\) of \(G\), the graph

\[
J=G-E(T)
\]

satisfies all three conditions:

\[
\boxed{
\begin{aligned}
\nu(J)&=k-1,\\
\tau(J)&=2(k-1),\\
J&\text{ has nonstandard triangle structure.}
\end{aligned}
}
\]

\paragraph{Proof.}

Minimality gives

\[
\tau(G)=2k+1
\]

and, for every edge \(e\),

\[
\nu(G-e)=k,\qquad \tau(G-e)=2k.
\]

For any triangle \(T\),

\[
\tau(G-E(T))
\geq\tau(G)-3
=2k-2.
\]

By minimality,

\[
\tau(G-E(T))
\leq2\nu(G-E(T)).
\]

Any packing of \(G-E(T)\) can be extended by adding \(T\), so

\[
\nu(G-E(T))\leq k-1.
\]

Equality follows throughout.

Now choose one edge \(e\) of \(T\), and set \(H=G-e\).

The graph \(H\) is obtained from \(J\) by adding two incident edges, and satisfies

\[
\nu(H)=\nu(J)+1,\qquad \tau(H)=2\nu(H).
\]

The incident-edge restoration theorem shows that \(J\) cannot have standard triangle structure.

\[
\square
\]

\section{Local consistency in a minimal counterexample}

Let $G$ be an edge-minimal counterexample, with $k=\nu(G)$ and $\tau(G)=2k+1$ as above. For every nonempty $F\subseteq E(G)$,
\[
 2k+1\le |F|+\tau(G-F)\le |F|+2\nu(G-F).
\]
Thus deleting $F$ costs at most $\lfloor (|F|-1)/2\rfloor$ triangles from a maximum packing. If $r$ edge-disjoint triangles and $s$ additional edges are specified, with $1\le r+s\le2$, delete the $3r+s$ edges involved. The bound gives a residual packing of size at least $k-r$. Restoring the $r$ specified triangles therefore gives a maximum packing containing those triangles and avoiding the $s$ additional edges. In particular, any two specified edges can be avoided, any specified triangle can be included while avoiding one other edge, and any two edge-disjoint triangles can be included together.

\section{Exact formulas for restoring a triangle}

Let $J$ be tight with $\nu(J)=m$, and let $a,b,c$ be independent vertices of $J$. Write $T=\{ab,ac,bc\}$ and $G=J+T$. When $J$ is obtained by deleting a triangle from an edge-minimal counterexample $G$, its partial restorations satisfy
\[
 \nu(J+S)=m+1,\qquad \tau(J+S)=2m+|S|
 \quad(\varnothing\ne S\subseteq T).
\]
For the cover number, adding $S$ to a minimum cover of $J$ gives the upper bound, and adding $T\setminus S$ to a cover of $J+S$ gives the lower bound from $\tau(G)=2m+3$. For proper nonempty $S$, minimality gives $2\nu(J+S)\ge2m+|S|$, and $\nu(J+S)\le m+1$. The full restoration follows by definition.

For a restored edge $uv\in T$, an \emph{extension wedge} consists of the two edges $ux,vx$ of $J$ for a common neighbor $x$. Call the type $uv$ active relative to a cover $C$ of $J$ if some wedge of this type avoids $C$. Write $A(C)\subseteq T$ for the active types. The exact covering formula is
\begin{equation}\label{eq:restoration-cover}
 \tau(G)=\min_{C\text{ covers }J}\bigl(|C|+\max\{1,|A(C)|\}\bigr).
\end{equation}
Indeed, once the old-edge portion $C$ of the cover is fixed, each active type requires its restored edge. Also, some restored edge must meet the triangle $abc$. These requirements suffice, because every other new triangle has one restored edge and one extension wedge. Since $\tau(J)=2m$, the condition $\tau(G)=2m+3$ is equivalent to both of the following.
\begin{enumerate}
\item Every size-$2m$ cover leaves all three types active.
\item Every size-$2m+1$ cover leaves at least two types active.
\end{enumerate}
Larger covers already cost at least $2m+3$ after meeting $abc$. We call these two conditions \emph{cover resistance}.

For the packing formula, let $\mathcal U$ run over sets of extension triangles using distinct restored edges and distinct outside vertices. Such triangles are edge-disjoint: for distinct types the only possible shared edge joins their shared distinguished vertex to a common outside vertex. Let $F(\mathcal U)$ be their union of old edges. Then
\begin{equation}\label{eq:restoration-pack}
 \nu(G)=\max\left\{m+1,\ \max_{\mathcal U}
 \bigl(|\mathcal U|+\nu(J-F(\mathcal U))\bigr)\right\}.
\end{equation}
The first branch uses $abc$. A packing leaving out $abc$ uses exactly a set $\mathcal U$ of extension triangles and a packing in the old graph avoiding their old edges. Thus $\nu(G)\ge m+2$ holds precisely when a pair of extension triangles can be added without losing any old packing triangles, or a triple can be added while losing at most one.

Tuza's conjecture is equivalent to the following statement: for every tight $J$ with three distinguished independent vertices satisfying cover resistance, at least one of these two packing augmentations exists. Indeed, if the conjecture holds, the exact cover number $2m+3$ forces a packing of size at least $m+2$. Conversely, deleting any triangle from an edge-minimal counterexample gives a tight $J$ satisfying cover resistance. Either augmentation would contradict $\nu(G)=m+1$. The augmentation statement remains open.

There are further restrictions in such a hypothetical counterexample. Restoring just one edge increases the packing number by one, so some wedge of each type is avoidable by a maximum packing of $J$. For two disjoint extension triangles, the packing loss in $J$ on deleting their four old edges is either one or two. It is at least one by~\eqref{eq:restoration-pack}. It is at most two because the residual cover number is at least $2m-4$ and the residual graph satisfies Tuza's inequality by edge-minimality. If the loss is two, local consistency in $G$ forces a compatible third extension triangle. A maximum packing containing the given pair would otherwise require $m-1$ triangles from the residual core, which has packing number $m-2$.

\subsection{A failed simultaneous avoidance rule}

Take $K_4$ on vertices $0,1,2,3$ and add pendant edges $04$ and $25$. The graph is tight with packing number one and cover number two. The open wedges $4$--$0$--$1$ and $5$--$2$--$3$ can each be avoided by some maximum packing. Avoiding both requires avoiding the opposite clique edges $01$ and $23$. Deleting these two edges leaves no triangle, so the two wedges cannot be avoided together by a maximum packing. The three-distinguished-vertex geometry in the augmentation statement is a genuine restriction.

\section{Exact small-graph checks}
\enlargethispage{\baselineskip}

The accompanying \texttt{verify-tight.py} computes packing and covering numbers by integer recurrences for all 1,253 NetworkX graph-atlas graphs. It recovers 44 tight graphs containing a triangle, whose triangle components are all $K_4$ or $K_5$, and verifies avoidance for all 250 open wedges in these graphs.

The verifier also checks all 167 independent triples in those tight graphs. For five triples, all three single-edge restorations increase both the cover and packing numbers by one. No triple has all three double-edge restorations increase the cover number by two, so none realizes the required full restoration profile. The finite checks include the simultaneous-avoidance counterexample above.

\section*{Disclosure}
This manuscript was prepared with AI assistance. Independent proof review and priority assessment are pending.

\begin{thebibliography}{9}
\bibitem{Tuza} Z. Tuza, \emph{A conjecture on triangles of graphs}, Graphs and Combinatorics 6 (1990), 373--380. \href{https://doi.org/10.1007/BF01787705}{Article.}
\end{thebibliography}
\end{document}
